\section{사용 하드웨어 --- 지원 하드웨어와 연결 핀 번호}

본 장은 \code{Pico/pins.py}와 \code{Pico/system\_config.py}에 정의된 실제 GPIO\footnote{\textbf{GPIO(General-Purpose
Input/Output, 범용 입출력)}: 마이크로컨트롤러의 디지털 신호 입출력 핀. 소프트웨어로 켜고
끄거나(출력) 외부 신호를 읽을(입력) 수 있다.}
핀 할당을 정리한다. 모든 핀은 라즈베리파이 피코(RP2040\footnote{\textbf{RP2040}: 라즈베리파이
재단이 설계한 마이크로컨트롤러 칩. 라즈베리파이 피코에 탑재되며 듀얼 코어 ARM Cortex-M0+
기반이다.})의 GP 번호 기준이다.

\subsection{GPIO 핀 할당 (\code{pins.py})}

\begin{longtable}{L{4.2cm} L{4.0cm} L{5.2cm}}
\toprule
\rowcolor{tablehead}\textbf{기능} & \textbf{상수명} & \textbf{GPIO} \\
\midrule
\endhead
UART 송신(블루투스) & \code{UART\_TX\_PIN} & GP0 \\
UART 수신(블루투스) & \code{UART\_RX\_PIN} & GP1 \\
DC 모터 방향 & \code{DCMOTOR\_DIR\_PIN} & GP8 \\
DC 모터 PWM & \code{DCMOTOR\_PWM\_PIN} & GP9 \\
I2C SDA(OLED) & \code{I2C\_SDA\_PIN} & GP4 \\
I2C SCL(OLED) & \code{I2C\_SCL\_PIN} & GP5 \\
초음파 트리거 & \code{ULTRASONIC\_TRIG\_PIN} & GP6 \\
초음파 에코 & \code{ULTRASONIC\_ECHO\_PIN} & GP7 \\
서보 휠(우) & \code{SERVO\_RIGHT\_PIN} & GP12 \\
서보 휠(좌) & \code{SERVO\_LEFT\_PIN} & GP13 \\
부저 & \code{BUZZER\_PIN} & GP15 \\
LED 배열(6개) & \code{LED\_PINS} & GP21, GP20, GP19, GP18, GP17, GP16 \\
BB탄 발사기 & \code{GUN\_PIN} & GP22 \\
자기장 센서 & \code{MAGSENSOR\_PIN} & GP26 \\
예비(미사용) & \code{RESERVED\_PINS} & GP10, GP11, GP14, GP27, GP28 \\
\bottomrule
\end{longtable}

\code{LED\_PINS}는 \code{[21, 20, 19, 18, 17, 16]} 순서로, LED0(메인)부터
LED5(오른쪽 끝)까지 6개이다. 명령에서 LED 번호는 0\textasciitilde 5로 지정한다.

핀과 주변장치의 물리적 연결을 한눈에 보면 그림~\ref{fig:wiring}과 같다.
RP2040의 GPIO는 인터페이스 유형에 따라 묶이는데, UART(GP0/1)는 BLE 통신,
I2C(GP4/5)는 OLED, 나머지는 PWM(서보\textbf{·}DC\textbf{·}LED\textbf{·}부저\textbf{·}발사기)과
디지털 I/O(초음파\textbf{·}자기)로 사용된다. GP10\textbf{·}11\textbf{·}14\textbf{·}27\textbf{·}28은
예비로 남겨 확장 모듈에 대비한다.

\begin{diagram}{RP2040 핀 \textbf{·} 주변장치 연결 연계도}{fig:wiring}
\node[dom, minimum height=58mm, text width=24mm] (mcu) {RP2040\\(Pico)};
% 왼쪽: 통신/센서
\node[hw, text width=30mm] (uart) at (-5.2,2.6) {BLE/UART 모듈\\{\scriptsize GP0 TX · GP1 RX}};
\node[hw, text width=30mm] (i2c) at (-5.2,1.1) {OLED (I2C)\\{\scriptsize GP4 SDA · GP5 SCL}};
\node[hw, text width=30mm] (us) at (-5.2,-0.4) {초음파 HC-SR04\\{\scriptsize GP6 TRIG · GP7 ECHO}};
\node[hw, text width=30mm] (dc) at (-5.2,-1.9) {DC 모터\\{\scriptsize GP8 DIR · GP9 PWM}};
% 오른쪽: 구동/출력
\node[hw, text width=30mm] (servo) at (5.2,2.6) {서보 휠\\{\scriptsize GP12 우 · GP13 좌}};
\node[hw, text width=30mm] (buz) at (5.2,1.3) {부저\\{\scriptsize GP15}};
\node[hw, text width=30mm] (led) at (5.2,0.0) {LED ×6\\{\scriptsize GP21·20·19·18·17·16}};
\node[hw, text width=30mm] (gun) at (5.2,-1.3) {발사기\\{\scriptsize GP22}};
\node[hw, text width=30mm] (mag) at (5.2,-2.6) {자기 센서\\{\scriptsize GP26}};
\foreach \n in {uart,i2c,us,dc} \draw[flow] (\n.east) -- (mcu.west);
\foreach \n in {servo,buz,led,gun,mag} \draw[flow] (mcu.east) -- (\n.west);
\node[font=\scriptsize\sffamily\color{rulegray}, below=1mm of mcu]
  {예비: GP10·11·14·27·28};
\end{diagram}

\subsection{지원 하드웨어 사양}

각 부품은 RP2040의 GPIO에 연결되어 PWM\footnote{\textbf{PWM(Pulse Width Modulation,
펄스 폭 변조)}: 일정 주기에서 신호가 켜져 있는 시간 비율(듀티)을 조절해 평균 전력을
바꾸는 기법. 모터 속도\textbf{·}LED 밝기 제어에 쓴다.}, I2C\footnote{\textbf{I2C}: 데이터선(SDA)과
클록선(SCL) 두 가닥으로 여러 장치를 연결하는 직렬 통신 버스. 본 프로젝트에서는 OLED 연결에
사용한다.}, ADC\footnote{\textbf{ADC(Analog-to-Digital Converter, 아날로그-디지털 변환기)}:
연속적인 아날로그 전압을 디지털 수치로 변환하는 회로.} 등 하드웨어 인터페이스로 제어된다.
모터처럼 회전 방향 전환이 필요한 경우 H-브리지\footnote{\textbf{H-브리지}: 네 개의 스위치로
모터에 흐르는 전류 방향을 바꿔 정\textbf{·}역 회전을 가능하게 하는 회로.} 회로를 사용한다.

\subsubsection{구동계}

\begin{longtable}{L{2.6cm} L{11.4cm}}
\toprule
\rowcolor{tablehead}\textbf{부품} & \textbf{사양 및 제어 방식} \\
\midrule
\endhead
서보 휠 & 연속회전 서보 2개. PWM 50\,Hz. 중립 듀티 약 5000(1.5\,ms),
  최소 1000\textasciitilde 최대 9000 범위. 좌우 보정 계수(\code{left\_calibration},
  \code{right\_calibration})로 직진성 교정. \\
DC 모터 & H-브리지 구동(GP8 방향, GP9 PWM). PWM 약 20\,kHz, 듀티 0\textasciitilde 65535.
  브레이크 변조 방식으로 양방향 대칭 토크 확보. \\
\bottomrule
\end{longtable}

\subsubsection{센서}

\begin{longtable}{L{2.6cm} L{11.4cm}}
\toprule
\rowcolor{tablehead}\textbf{부품} & \textbf{사양 및 측정 방식} \\
\midrule
\endhead
초음파 거리 & HC-SR04. TRIG(GP6) 펄스 후 ECHO(GP7) 에코 시간 측정,
  \code{distance\_cm = (duration/2) * 0.0343}로 환산. \\
자기 센서 & 디지털 입력(GP26), 풀업. 자석 근접 시 0(액티브 로우)을 1로 변환하여 검출. \\
온도 & RP2040 내장 온도 센서(ADC4). 상태 보고용. \\
\bottomrule
\end{longtable}

\subsubsection{출력 \textbf{·} 표시 \textbf{·} 음향}

\begin{longtable}{L{2.6cm} L{11.4cm}}
\toprule
\rowcolor{tablehead}\textbf{부품} & \textbf{사양 및 제어 방식} \\
\midrule
\endhead
LED 배열 & 6개 PWM 출력(약 20\,kHz). 밝기 0.0\textasciitilde 1.0을 듀티
  0\textasciitilde 65535로 매핑. 스와이프 애니메이션 지원. \\
부저 & PWM 톤(주파수 가변). 비차단 재생: 명령은 즉시 반환되고 메인 루프의
  \code{buzzer.update()}가 지속시간 경과 후 자동 정지. \\
OLED & SSD1306, 128$\times$64 모노크롬, I2C(0) 400\,kHz, 주소 0x3c.
  텍스트/사각형/아이콘(로버\textbf{·}화성\textbf{·}눈 뜬/감은) 출력. \\
발사기 & BB탄 발사 모터(GP22). PWM 약 1\,kHz. 소프트 스타트 후 단발 발사
  (\code{fire\_once()}). \\
\bottomrule
\end{longtable}

\subsection{무선 통신 모듈}

피코는 UART(GP0/GP1)를 통해 BLE UART 브리지 모듈(HM-10 또는 BT05,
CC2541 계열 Bluetooth 4.0)과 연결된다. 통신 파라미터는 다음과 같다.

\begin{itemize}
  \item UART: \code{UART(0)}, \textbf{9600 baud}, 8N1
  \item 프레임 종단: 개행 문자(\code{\textbackslash n})
  \item BLE GATT\footnote{\textbf{GATT(Generic Attribute Profile)}: BLE에서 데이터를 서비스와
        특성(characteristic) 단위로 구조화해 주고받도록 정의한 프로토콜.} UART 서비스
        UUID\footnote{\textbf{UUID(범용 고유 식별자)}: 중복 가능성이 거의 없는 128비트 식별자.
        BLE 서비스\textbf{·}특성을 구분하는 데 쓴다.}: \code{0000ffe0-0000-1000-8000-00805f9b34fb}
  \item BLE 특성(characteristic) UUID: \code{0000ffe1-0000-1000-8000-00805f9b34fb}
  \item 장치 이름 최대 길이: \textbf{10자} (BT05/HMSoft 제약)
\end{itemize}

\subsection{핀 \textbf{·} 모듈의 런타임 재구성}

핀 할당과 모듈 활성 상태는 펌웨어 재기록 없이 변경할 수 있다.

\begin{itemize}
  \item \code{SET\_PIN,핀이름,번호} --- 핀 번호를 \code{system.json}에 저장.
        유효 범위 0\textasciitilde 28. 다음 부팅 시 적용.
  \item \code{SET\_MODULE,모듈명,0|1} --- 모듈 활성/비활성 토글.
  \item \code{GET\_MODULES} --- 전체 모듈 활성 상태 조회.
\end{itemize}

\begin{infobox}[title=설계 의도]
핀과 모듈을 \code{pins.py}와 \code{system.json}에 중앙 집중화함으로써, 교실
현장에서 부품 구성이 다른 로버(예: 발사기 미장착)도 동일한 펌웨어로 운용할 수
있다. 운영자는 코드를 고치지 않고 설정만으로 하드웨어 구성을 맞춘다.
\end{infobox}

\begin{teachbox}{핀과 인터페이스의 중앙 관리}
\teachhead{설계 의도}
\begin{itemize}
  \item 핀 번호를 \code{pins.py} 한 곳에 모으고 런타임 재구성(\code{SET\_PIN})을 지원해, 코드를 고치지 않고 하드웨어 구성을 바꾼다.
  \item PWM\textbf{·}I2C\textbf{·}ADC 같은 인터페이스를 부품 성격에 맞게 골라 썼다.
\end{itemize}
\teachhead{분석할 때 핵심 질문}
\begin{itemize}
  \item 핀 번호가 코드 곳곳에 흩어져 있었다면 어떤 불편이 생길까?
  \item 모터엔 PWM, OLED엔 I2C, 온도엔 ADC를 쓴 \emph{이유}는 무엇인가?
\end{itemize}
\teachhead{점검 요소}
\begin{itemize}
  \item LED 핀 하나를 예비 핀으로 옮기고(설정 변경 후 재부팅) 정상 동작을 확인하라.
\end{itemize}
\end{teachbox}

